\chapter{Complexity — Oblivious}
%%% generated by the blueprint pipeline from TCSlib.Complexity.TuringMachine.Robustness.Oblivious
\section{Overview}

This module proves the correctness of the oblivious simulation (Arora--Barak,
Remark~1.7 and Exercise~1.5): it establishes the data machine's exact sweep
identities, the preparation invariant carrying the copied input through
initialization, the macrostep-by-macrostep tracking of the simulated source
run, and finally the theorem that every language decidable in a
time-constructible bound $T$ is decided by an oblivious machine within
$c\,(T(n)+1)^2$ steps.

\section{Declarations}

\begin{definition}[Logical data simulator]
\lean{Complexity.dataTM}
\label{Complexity.dataTM}
\leanok
\uses{Complexity.OblSymbol, Complexity.decorateTM, Complexity.obliviousDataInit, Complexity.obliviousSchedule, Complexity.obliviousVisit, Turing.FinTM}
Given binary machines $W$ (the clock witness) and $M$ (the source, with $k$ work
tapes) and a padding constant $a$, the logical simulator is the machine over the
extended simulation alphabet obtained by decorating the length-only schedule
controller built from $W$ and $a$ with $k+1$ additional data tapes tracked at the
guide-tape head, starting from the initial data state of $M$ and driven by the
data transducer of $W$, $M$, and $a$.
\end{definition}

\begin{definition}[Macrostep configuration with data]
\lean{Complexity.dataCfg}
\label{Complexity.dataCfg}
\leanok
\uses{Complexity.OblData, Complexity.OblPhase, Complexity.OblSymbol, Complexity.decoratedCfg, Complexity.macroCfg, Turing.Cfg, Turing.FinTM}
Let $W$ and $M$ be binary machines and $a$ a padding constant.  Given a base
schedule configuration, a schedule phase $q$, unary and guide head positions
$u, g \in \bbz$, a finite data register $d$, contents for the $k+1$ data tapes,
and an output list, the associated configuration of the logical simulator is the
decorated form of the macrostep configuration of the base at phase $q$ with
heads $u$ and $g$, carrying $d$, the data tapes, and the output.
\end{definition}

\begin{lemma}[Exact form of one simulator step]
\lean{Complexity.dataCfg_step}
\label{Complexity.dataCfg_step}
\leanok
\uses{Complexity.OblData, Complexity.OblPhase, Complexity.OblSymbol, Complexity.dataCfg, Complexity.dataTM, Complexity.decoratedCfg, Complexity.decoratedCfg_step, Complexity.macroCfg, Complexity.obliviousDataInit, Complexity.obliviousSchedule, Complexity.obliviousVisit, Complexity.setupWrite, Turing.Cfg, Turing.Cfg.inputSymbol, Turing.Cfg.workTapeSymbols, Turing.FinTM, Turing.MultiTapeTM.step}
\difficulty{2}
Let $W$ and $M$ be binary machines and $a$ a padding constant.  One step of the
logical simulator on a macrostep configuration with data equals the decorated
form of one schedule-controller step, where the new finite data register, the
single cell written under the guide head on each data tape, and the appended
emission are exactly those returned by the data transducer evaluated at the
current phase, register, input read, schedule tape reads, and the data cells
under the guide head.
\end{lemma}

\begin{lemma}[Forward visit on the data tapes]
\lean{Complexity.dataCfg_forward_step}
\label{Complexity.dataCfg_forward_step}
\leanok
\uses{Complexity.OblPayload, Complexity.OblPhase, Complexity.OblSymbol, Complexity.dataCell, Complexity.dataCfg, Complexity.dataCfg_step, Complexity.dataTM, Complexity.guideTape, Complexity.guideTape_right, Complexity.macroCfg_forward, Complexity.macroCfg_guide, Complexity.obliviousVisit, Complexity.setupWrite, Turing.Cfg, Turing.FinTM, Turing.MultiTapeTM.step}
\difficulty{3}
Suppose the base's guide tape is the fixed guide of radius $R$ and the simulator
is in the forward phase with guide head at $g \neq R + 1$.  One step keeps the
saved state, answer bit, and reads, moves the guide head to $g+1$, replaces each
neighbor register by the current payload decoded at $g$, and rewrites cell $g$
of each data tape to a data cell carrying its unchanged current payload with the
old neighbor register cached; no output is emitted.
\end{lemma}

\begin{lemma}[Forward boundary turn]
\lean{Complexity.dataCfg_forward_turn}
\label{Complexity.dataCfg_forward_turn}
\leanok
\uses{Complexity.OblPayload, Complexity.OblPhase, Complexity.OblSymbol, Complexity.blankPayload, Complexity.dataCfg, Complexity.dataCfg_step, Complexity.dataTM, Complexity.guideTape, Complexity.guideTape_right, Complexity.macroCfg_forward, Complexity.macroCfg_guide, Complexity.obliviousVisit, Complexity.setupWrite, Turing.Cfg, Turing.FinTM, Turing.MultiTapeTM.step}
\difficulty{3}
Suppose the base's guide tape is the fixed guide of radius $R$.  From the
forward phase with guide head at $R + 1$, one simulator step enters the backward
phase at position $R$, resets every neighbor register to the blank payload, and
leaves the saved state, answer bit, reads, data tapes, and output unchanged.
\end{lemma}

\begin{lemma}[Backward visit on the data tapes]
\lean{Complexity.dataCfg_backward_step}
\label{Complexity.dataCfg_backward_step}
\leanok
\uses{Complexity.OblPayload, Complexity.OblPhase, Complexity.OblSymbol, Complexity.blankPayload, Complexity.dataCell, Complexity.dataCfg, Complexity.dataCfg_step, Complexity.dataTM, Complexity.guideTape, Complexity.guideTape_left, Complexity.macroCfg_backward, Complexity.macroCfg_guide, Complexity.obliviousSourceMove, Complexity.obliviousVisit, Complexity.setupWrite, Turing.Cfg, Turing.FinTM, Turing.MultiTapeTM.step}
\difficulty{3}
Suppose the base's guide tape is the fixed guide of radius $R$ and the simulator
is in the backward phase with guide head at $g \neq -R - 1$.  One step moves the
guide head to $g - 1$, replaces each neighbor register by the current payload
decoded at $g$, and rewrites cell $g$ of each data lane to a data cell with
blank cached neighbor whose current payload is chosen by that lane's virtual
head direction: the cached left neighbor for a leftward move, the unchanged
current payload for a stationary move, and the carried register for a rightward
move.
\end{lemma}

\begin{lemma}[Backward boundary turn]
\lean{Complexity.dataCfg_backward_turn}
\label{Complexity.dataCfg_backward_turn}
\leanok
\uses{Complexity.OblData, Complexity.OblPhase, Complexity.OblSymbol, Complexity.dataCfg, Complexity.dataCfg_step, Complexity.dataTM, Complexity.guideTape, Complexity.guideTape_left, Complexity.macroCfg_backward, Complexity.macroCfg_guide, Complexity.obliviousVisit, Complexity.setupWrite, Turing.Cfg, Turing.FinTM, Turing.MultiTapeTM.step}
\difficulty{3}
Suppose the base's guide tape is the fixed guide of radius $R$.  From the
backward phase with guide head at $-R - 1$, one simulator step enters the return
phase at position $-R$ and leaves the finite data register, data tapes, and
output unchanged.
\end{lemma}

\begin{lemma}[Invariant of the forward scan]
\lean{Complexity.dataCfg_forward_prefix}
\label{Complexity.dataCfg_forward_prefix}
\leanok
\uses{Complexity.OblPayload, Complexity.OblPhase, Complexity.OblSymbol, Complexity.dataCell, Complexity.dataCfg, Complexity.dataCfg_forward_step, Complexity.dataTM, Complexity.guideTape, Turing.Cfg, Turing.FinTM, Turing.MultiTapeTM.runFrom, Turing.MultiTapeTM.runFrom_succ_eq_step', Turing.MultiTapeTM.runFrom_zero}
\difficulty{5}
Suppose the base's guide tape is the fixed guide of radius $R$, let $f$ assign a
payload vector to every integer coordinate, and suppose every current payload of
the data tapes decodes to $f$.  For every $n \le 2R + 1$, running the simulator
$n$ steps from the forward phase at $-R$ with neighbor registers $f(-R-1)$ stays
in the forward phase, reaching position $-R+n$ with neighbor registers
$f(-R+n-1)$, on tapes whose current payloads still decode to $f$ and whose
cached left-neighbor payloads equal $f(z-1)$ at every visited coordinate
$-R \le z < -R + n$.
\end{lemma}

\begin{lemma}[Invariant of the backward scan]
\lean{Complexity.dataCfg_backward_prefix}
\label{Complexity.dataCfg_backward_prefix}
\leanok
\uses{Complexity.OblPayload, Complexity.OblPhase, Complexity.OblSymbol, Complexity.blankPayload, Complexity.dataCell, Complexity.dataCfg, Complexity.dataCfg_backward_step, Complexity.dataTM, Complexity.guideTape, Complexity.obliviousSourceMove, Turing.Cfg, Turing.FinTM, Turing.MultiTapeTM.runFrom, Turing.MultiTapeTM.runFrom_succ_eq_step', Turing.MultiTapeTM.runFrom_zero}
\difficulty{5}
Suppose the base's guide tape is the fixed guide of radius $R$, let $f$ assign a
payload vector to every coordinate, and suppose the current payloads of the data
tapes decode to $f$ while the cached left neighbors equal $f(z-1)$ throughout
$[-R, R]$.  For every $n \le 2R + 1$, running the simulator $n$ steps from the
backward phase at $R$ with neighbor registers $f(R+1)$ stays in the backward
phase, reaching position $R - n$ with neighbor registers $f(R-n+1)$, on tapes
whose current payload on lane $i$ at coordinate $z$ is the shifted value
$f(z + m_i)$ for visited coordinates $R - n < z \le R$ — where $m_i$ is the
virtual head direction of lane $i$ selected from the saved state and reads —
and the unshifted value $f(z)$ elsewhere, while the cached left neighbors remain
$f(z-1)$ for $-R \le z \le R - n$.
\end{lemma}

\begin{lemma}[Exact outward positioning scan]
\lean{Complexity.dataCfg_seek_run}
\label{Complexity.dataCfg_seek_run}
\leanok
\uses{Complexity.OblPayload, Complexity.OblPhase, Complexity.OblSymbol, Complexity.blankPayload, Complexity.dataCfg, Complexity.dataCfg_step, Complexity.dataTM, Complexity.guideTape, Complexity.guideTape_left, Complexity.macroCfg_guide, Complexity.macroCfg_seek, Complexity.obliviousVisit, Complexity.setupWrite, Turing.Cfg, Turing.FinTM, Turing.MultiTapeTM.runFrom, Turing.MultiTapeTM.runFrom_succ_eq_step', Turing.MultiTapeTM.runFrom_zero}
\difficulty{4}
Suppose the base's guide tape is the fixed guide of radius $R$.  Running the
simulator exactly $R + 2$ steps from the leftward-seek phase with guide head at
the origin and blank neighbor registers reaches the forward phase at $-R$ with
blank neighbor registers, leaving the saved state, answer bit, reads, data
tapes, and output unchanged.
\end{lemma}

\begin{lemma}[Exact return to the origin]
\lean{Complexity.dataCfg_center_run}
\label{Complexity.dataCfg_center_run}
\leanok
\uses{Complexity.OblPayload, Complexity.OblPhase, Complexity.OblSymbol, Complexity.dataCfg, Complexity.dataCfg_step, Complexity.dataTM, Complexity.guideTape, Complexity.guideTape_origin, Complexity.macroCfg_center, Complexity.macroCfg_guide, Complexity.obliviousSourceAction, Complexity.obliviousVisit, Complexity.setupWrite, Turing.Cfg, Turing.FinTM, Turing.MultiTapeTM.runFrom, Turing.MultiTapeTM.runFrom_succ_eq_step', Turing.MultiTapeTM.runFrom_zero}
\difficulty{4}
Suppose the base's guide tape is the fixed guide of radius $R$.  Running the
simulator exactly $R + 1$ steps from the return phase with guide head at $-R$
and unary head at $u$ reaches the counter-test phase with unary head $u + 1$ and
guide head at the origin; the simulated state is replaced by the successor state
of the source action selected from the saved state and reads, while the answer
bit, reads, neighbor registers, data tapes, and output are unchanged.
\end{lemma}

\begin{lemma}[The counter test starts a sweep]
\lean{Complexity.dataCfg_check}
\label{Complexity.dataCfg_check}
\leanok
\uses{Complexity.OblData, Complexity.OblPhase, Complexity.OblSymbol, Complexity.blankPayload, Complexity.dataCell, Complexity.dataCfg, Complexity.dataCfg_step, Complexity.dataTM, Complexity.macroCfg_check, Complexity.macroCfg_unary, Complexity.obliviousSourceAction, Complexity.obliviousVisit, Complexity.setupWrite, Turing.Cfg, Turing.FinTM, Turing.MultiTapeTM.step}
\difficulty{3}
Suppose the base's unary tape holds a unit marker at $u$.  One simulator step
from the counter-test phase saves the current payloads decoded at the origin as
the new reads, updates the answer register to the bit emitted by the source
action selected from the saved state and those reads (keeping the old answer
when nothing is emitted), performs that action's optional work-lane writes at
the origin, clears the neighbor registers to blank, and enters the leftward-seek
phase; no output is emitted.
\end{lemma}

\begin{lemma}[Origin writes decode to the source's writes]
\lean{Complexity.dataCfg_written}
\label{Complexity.dataCfg_written}
\leanok
\uses{Complexity.OblSymbol, Complexity.blankPayload, Complexity.dataCell, Complexity.setupWrite, Complexity.sourcePayload, Complexity.sourceTotalAction, Complexity.sourceWrittenPayload, Turing.Cfg, Turing.FinTM}
\difficulty{4}
Let $M$ be a binary machine and $c$ a configuration of $M$, and suppose the
current payloads of the data tapes decode to the head-centered payload vectors
of $c$.  After writing at the origin of each work lane the symbol optionally
written by the totalized action of $M$ at $c$ (with blank cached neighbor, the
input lane untouched), the current payloads decode to the written payload
vectors of $c$ under that action.
\end{lemma}

\begin{lemma}[One full pair of data sweeps]
\lean{Complexity.dataCfg_sweeps}
\label{Complexity.dataCfg_sweeps}
\leanok
\uses{Complexity.OblPayload, Complexity.OblPhase, Complexity.OblSymbol, Complexity.blankPayload, Complexity.dataCell, Complexity.dataCfg, Complexity.dataCfg_backward_prefix, Complexity.dataCfg_backward_turn, Complexity.dataCfg_center_run, Complexity.dataCfg_forward_prefix, Complexity.dataCfg_forward_turn, Complexity.dataCfg_seek_run, Complexity.dataTM, Complexity.guideTape, Complexity.obliviousSourceAction, Complexity.obliviousSourceMove, Turing.Cfg, Turing.FinTM, Turing.MultiTapeTM.runFrom, Turing.MultiTapeTM.runFrom_add, Turing.MultiTapeTM.runFrom_succ_eq_step}
\difficulty{5}
Suppose the base's guide tape is the fixed guide of radius $R$, the current
payloads of the data tapes decode to $f$, and the boundary values $f(-R-1)$ and
$f(R+1)$ are blank.  Then running the simulator exactly $6R + 7$ steps from the
leftward-seek phase at the origin with unary head $u$ and blank neighbor
registers reaches the counter-test phase with unary head $u + 1$, simulated
state equal to the successor state of the source action selected from the saved
state and reads, and neighbor registers $f(-R)$, on tapes whose current payload
on lane $i$ at coordinate $z$ is $f(z + m_i)$ for $-R \le z \le R$ and $f(z)$
outside, where $m_i$ is that lane's selected virtual head direction.
\end{lemma}

\begin{lemma}[Origin writes are local]
\lean{Complexity.sourceWrittenPayload_away}
\label{Complexity.sourceWrittenPayload_away}
\leanok
\uses{Complexity.sourcePayload, Complexity.sourceTotalAction, Complexity.sourceWrittenPayload, Turing.Cfg, Turing.FinTM}
\difficulty{2}
For a binary machine $M$, a configuration $c$ of $M$, and any offset $z \neq 0$,
the written payload vector of $c$ under the totalized action of $M$ agrees at
$z$ with the head-centered payload vector of $c$.
\end{lemma}

\begin{lemma}[One macrostep simulates one source step]
\lean{Complexity.dataCfg_simulates}
\label{Complexity.dataCfg_simulates}
\leanok
\uses{Complexity.OblPayload, Complexity.OblPhase, Complexity.OblSymbol, Complexity.blankPayload, Complexity.dataCell, Complexity.dataCfg, Complexity.dataCfg_check, Complexity.dataCfg_sweeps, Complexity.dataCfg_written, Complexity.dataTM, Complexity.guideTape, Complexity.lastOutput_append, Complexity.obliviousSourceAction_correct, Complexity.obliviousSourceMove_correct, Complexity.setupWrite, Complexity.sourcePayload, Complexity.sourcePayload_apply, Complexity.sourceTotalAction, Complexity.sourceTotalAction_apply, Complexity.sourceWrittenPayload, Complexity.sourceWrittenPayload_away, Turing.Cfg, Turing.FinTM, Turing.MultiTapeTM.runFrom, Turing.MultiTapeTM.runFrom_succ_eq_step, Turing.MultiTapeTM.step}
\difficulty{5}
Suppose the base's guide tape is the fixed guide of radius $R$ and its unary
tape holds a unit marker at $u$, and let $c$ be a configuration of the source
$M$ whose head-centered payload vectors are decoded by the data tapes and are
blank outside $[-R, R]$, as are those of its one-step successor.  Then exactly
$6R + 8$ simulator steps from the counter-test phase carrying the state and
last-emitted bit of $c$ reach the counter-test phase at unary head $u + 1$
carrying the state and last-emitted bit of the successor of $c$, with data tapes
decoding the successor's head-centered payload vectors and an unchanged output
stream.
\end{lemma}

\begin{definition}[Payloads of a partial input copy]
\lean{Complexity.copiedPayload}
\label{Complexity.copiedPayload}
\leanok
\uses{Complexity.OblPayload, Complexity.blankPayload, Complexity.inputPayload, Turing.FinTM}
For a binary machine $M$ with $k$ work tapes, a word $x$, and a count $j$, the
copied payload at coordinate $z$ is blank on every work lane; on the input lane
it is the virtual input payload of $x$ at $z$ when $z = -1$ (the copied left
boundary) or $0 \le z < j$, and blank otherwise.
\end{definition}

\begin{lemma}[The full copy represents the initial configuration]
\lean{Complexity.copiedPayload_full}
\label{Complexity.copiedPayload_full}
\leanok
\uses{Complexity.blankPayload, Complexity.copiedPayload, Complexity.inputPayload_outside, Complexity.sourcePayload, Turing.Cfg.init, Turing.FinTM, Turing.MultiTapeTM.initCfg}
\difficulty{3}
For every binary machine $M$ and word $x$, the copied payloads of count
$|x| + 1$ coincide with the head-centered payload vectors of the initial
configuration of $M$ on $x$: the input lane carries the complete virtual input
including both boundary tags, and every work lane is blank.
\end{lemma}

\begin{definition}[Input-lane copy write]
\lean{Complexity.copyDataWrite}
\label{Complexity.copyDataWrite}
\leanok
\uses{Complexity.OblPayload, Complexity.OblSymbol, Complexity.blankPayload, Complexity.setupWrite, Turing.FinTM}
For a binary machine $M$, given contents of the data tapes, a coordinate $z$,
and a payload $p$, the copy write updates only the input lane, placing at $z$ a
data cell with current payload $p$ and blank cached neighbor; the work lanes are
unchanged.
\end{definition}

\begin{lemma}[Copying the left boundary]
\lean{Complexity.copyDataWrite_left}
\label{Complexity.copyDataWrite_left}
\leanok
\uses{Complexity.OblSymbol, Complexity.blankPayload, Complexity.copiedPayload, Complexity.copyDataWrite, Complexity.dataCell, Complexity.inputPayload, Complexity.setupWrite, Turing.FinTM, Turing.FinTM.bufferTape}
\difficulty{3}
If every current payload of the data tapes is blank, then after the copy write
placing the left-boundary payload (blank content with tag $1$) at coordinate
$-1$, the current payloads are exactly the copied payloads of count $0$.
\end{lemma}

\begin{lemma}[Copying one more input cell]
\lean{Complexity.copyDataWrite_next}
\label{Complexity.copyDataWrite_next}
\leanok
\uses{Complexity.OblSymbol, Complexity.copiedPayload, Complexity.copyDataWrite, Complexity.dataCell, Complexity.inputPayload, Complexity.setupWrite, Turing.FinTM}
\difficulty{3}
If the current payloads of the data tapes are the copied payloads of count $j$,
then after the copy write placing the virtual input payload of $x$ at position
$j$ at coordinate $j$, the current payloads are the copied payloads of count
$j + 1$.
\end{lemma}

\begin{definition}[Initial data assertions]
\lean{Complexity.initialContent}
\label{Complexity.initialContent}
\leanok
\uses{Complexity.OblData, Complexity.OblPayload, Complexity.OblSymbol, Complexity.dataCell, Complexity.obliviousDataInit, Turing.FinTM}
For a binary machine $M$, a payload assignment $f$, a finite data register, data
tapes, and an output list, the initial-content predicate states that the
register is the initial data state of $M$, the output is empty, and the current
payload decoded from the data tapes at every lane and coordinate equals the
value of $f$ there.
\end{definition}

\begin{definition}[Phase-indexed initialization invariant]
\lean{Complexity.prepCondition}
\label{Complexity.prepCondition}
\leanok
\uses{Complexity.OblData, Complexity.OblPhase, Complexity.OblSymbol, Complexity.blankPayload, Complexity.copiedPayload, Complexity.initialContent, Turing.FinTM}
For a source machine $M$, a word $x$, an optional schedule phase, an input head
position, unary and guide head positions, a unary tape, a data register, data
tapes, and an output list, the preparation condition records the information
preserved during initialization: except in the very first phase, the unary tape
carries its origin sentinel exactly at coordinate $0$; the data register and
output are still initial; the data tapes hold blank payloads before copying, the
copied payloads of the appropriate count during copying, and the full input copy
in the later allocation phases; and the input, unary, and guide heads sit at
their phase-prescribed positions.  During the macrostep cycle only the lower
bound $1 \le u$ on the unary head is kept, and at the counter-test phase the
full copy is required exactly when $u = 1$; a halted schedule imposes no
condition.
\end{definition}

\begin{definition}[Configuration form of the initialization invariant]
\lean{Complexity.prepInvariant}
\label{Complexity.prepInvariant}
\leanok
\uses{Complexity.OblData, Complexity.OblPhase, Complexity.OblSymbol, Complexity.oblEmbed, Complexity.prepCondition, Turing.Cfg, Turing.FinTM}
For a schedule configuration $c$ on the embedded input, the preparation
invariant is the preparation condition instantiated at the state of $c$, its
input head position, its unary and guide head positions, and its unary tape
content, together with the given data register, data tapes, and output.
\end{definition}

\begin{lemma}[Blank input reads occur only at the boundaries]
\lean{Complexity.inputSymbol_blank}
\label{Complexity.inputSymbol_blank}
\leanok
\uses{Turing.Cfg, Turing.Cfg.inputSymbol}
\difficulty{2}
For any configuration of a multi-tape machine on an input word $x$, the input
read is blank if and only if the input head is at position $0$ or at position
$|x| + 1$.
\end{lemma}

\begin{lemma}[Input read at the copy position]
\lean{Complexity.copy_input_read}
\label{Complexity.copy_input_read}
\leanok
\uses{Complexity.OblSymbol, Complexity.inputSymbol_blank, Complexity.oblEmbed, Turing.Cfg, Turing.Cfg.inputSymbol, Turing.inputSymbolInner}
\difficulty{3}
Let $x$ be a binary word embedded letterwise into the simulation alphabet, and
let a configuration on this embedded input have its input head at position
$j + 1$ with $j \le |x|$.  Then the input read is the bit symbol of $x_j$ when
$j < |x|$, and blank when $j = |x|$; this includes the empty input.
\end{lemma}

\begin{lemma}[Copied reads give the virtual input payload]
\lean{Complexity.copy_input_payload}
\label{Complexity.copy_input_payload}
\leanok
\uses{Complexity.OblSymbol, Complexity.clockBit, Complexity.copy_input_read, Complexity.inputPayload, Complexity.oblEmbed, Turing.Cfg, Turing.Cfg.inputSymbol, Turing.FinTM.bufferTape}
\difficulty{3}
Let $x$ be a binary word embedded into the simulation alphabet and let a
configuration on the embedded input have its input head at position $j + 1$ with
$j \le |x|$.  Then the payload whose content is the clock-bit decoding of the
input read and whose tag is the right-boundary tag $2$ when the read is blank
and the ordinary tag $0$ otherwise is exactly the virtual input payload of $x$
at position $j$.
\end{lemma}

\begin{lemma}[The invariant after a schedule action]
\lean{Complexity.prepInvariant_action}
\label{Complexity.prepInvariant_action}
\leanok
\uses{Complexity.OblData, Complexity.OblPhase, Complexity.OblSymbol, Complexity.oblAction, Complexity.oblEmbed, Complexity.prepCondition, Complexity.prepInvariant, Complexity.setupWrite, Turing.Action.apply, Turing.Cfg, Turing.FinTM, Turing.moveInputPos}
\difficulty{2}
Applying to a schedule configuration a non-clock action of the schedule
controller — with a target phase, an input move, and optional writes and moves
on the three distinguished tapes — the preparation invariant of the resulting
configuration equals the preparation condition at the target phase, the moved
input position, the unary and guide head positions shifted by the action's
moves, and the unary tape after the action's optional write, with the data
register, data tapes, and output unchanged.
\end{lemma}

\begin{lemma}[The origin survives unit writes]
\lean{Complexity.unaryOrigin_update}
\label{Complexity.unaryOrigin_update}
\leanok
\uses{Complexity.OblSymbol}
\difficulty{2}
If a tape carries the origin sentinel exactly at coordinate $0$, then after
writing a unit marker at any coordinate $u \ge 1$ it still carries the origin
sentinel exactly at coordinate $0$.
\end{lemma}

\begin{lemma}[Preparation is invariant under transitions]
\lean{Complexity.prepInvariant_step}
\label{Complexity.prepInvariant_step}
\leanok
\uses{Complexity.OblData, Complexity.OblPhase, Complexity.OblSymbol, Complexity.blankPayload, Complexity.clockBit, Complexity.copyDataWrite, Complexity.copyDataWrite_left, Complexity.copyDataWrite_next, Complexity.copy_input_payload, Complexity.copy_input_read, Complexity.initialContent, Complexity.inputSymbol_blank, Complexity.oblEmbed, Complexity.obliviousSchedule, Complexity.obliviousVisit, Complexity.prepCondition, Complexity.prepInvariant, Complexity.prepInvariant_action, Complexity.setupWrite, Complexity.unaryOrigin_update, Turing.Action.apply, Turing.Cfg, Turing.Cfg.inputSymbol, Turing.Cfg.workTapeSymbols, Turing.FinTM, Turing.FinTM.moveInputPos_neg_val, Turing.MultiTapeTM.step, Turing.moveInputPos, Turing.moveInputPos_pos_of_ne_right}
\difficulty{6}
Let a schedule configuration $c$ on the embedded input be in a live phase and
satisfy the preparation invariant with a given data register, data tapes, and
output.  Then the preparation invariant holds for the configuration after one
step of the schedule controller, taken with the register, guide-position tape
writes, and appended emission produced by the data transducer's visit at $c$.
\end{lemma}

\begin{lemma}[The invariant holds at all times]
\lean{Complexity.prepInvariant_run}
\label{Complexity.prepInvariant_run}
\leanok
\uses{Complexity.dataCell, Complexity.dataTM, Complexity.decoratedCfg, Complexity.decoratedCfg_init, Complexity.decoratedCfg_step, Complexity.initialContent, Complexity.oblEmbed, Complexity.obliviousDataInit, Complexity.obliviousSchedule, Complexity.obliviousVisit, Complexity.prepCondition, Complexity.prepInvariant, Complexity.prepInvariant_step, Complexity.setupWrite, Turing.Cfg.inputSymbol, Turing.Cfg.workTapeSymbols, Turing.FinTM, Turing.MultiTapeTM.initCfg, Turing.MultiTapeTM.runFrom, Turing.MultiTapeTM.runFrom_succ_eq_step', Turing.MultiTapeTM.step, Turing.MultiTapeTM.step_of_halt}
\difficulty{5}
For all binary machines $W$ and $M$, padding constant $a$, word $x$, and time
$t$, there exist a data register, data tapes, and an output list such that the
$t$-step run of the logical simulator from its initial configuration on the
embedded input equals the decorated form of the schedule controller's $t$-step
run with these components, and the preparation invariant holds for that schedule
configuration with these components.
\end{lemma}

\begin{lemma}[Reassembling a configuration from its own heads]
\lean{Complexity.macroCfg_self}
\label{Complexity.macroCfg_self}
\leanok
\uses{Complexity.OblPhase, Complexity.OblSymbol, Complexity.finThree_cases, Complexity.macroCfg, Turing.Cfg, Turing.FinTM}
\difficulty{3}
Rebuilding a schedule configuration in macrostep form from its own state and its
actual unary and guide head positions returns the configuration itself.
\end{lemma}

\begin{lemma}[The schedule reaches the first macrostep]
\lean{Complexity.obliviousSchedule_ready}
\label{Complexity.obliviousSchedule_ready}
\leanok
\uses{Complexity.budgetValue, Complexity.budgetValue_bits, Complexity.clockStageCfg, Complexity.clockStageCfg_captures, Complexity.clockStageCfg_setup, Complexity.guideTape, Complexity.oblEmbed, Complexity.obliviousSchedule, Complexity.setupCfg, Complexity.setupCfg_initializes, Complexity.unaryTape, Turing.FinTM, Turing.FinTM.ComputesInTime, Turing.MultiTapeTM.initCfg, Turing.MultiTapeTM.runFrom, Turing.MultiTapeTM.runFrom_add}
\difficulty{5}
Let $W$ compute on every input $x$ the binary expansion of $T(|x|)$ within
$b\,(T(|x|)+1)$ steps, and suppose $n \le T(n)$ for all $n$.  Then for every
word $x$ there is a time $t$ at which the schedule controller built from $W$ and
$a$, run on the embedded input, is in the counter-test phase with unary head at
$1$, guide head at the origin, unary counter tape with $(a+1)(T(|x|)+1)$
markers, and guide tape of radius $3(a+1)(T(|x|)+1)$.
\end{lemma}

\begin{lemma}[The simulator reaches the first macrostep]
\lean{Complexity.dataTM_ready}
\label{Complexity.dataTM_ready}
\leanok
\uses{Complexity.OblPhase, Complexity.OblSymbol, Complexity.copiedPayload_full, Complexity.dataCell, Complexity.dataCfg, Complexity.dataTM, Complexity.guideTape, Complexity.macroCfg_self, Complexity.oblEmbed, Complexity.obliviousDataInit, Complexity.obliviousSchedule, Complexity.obliviousSchedule_ready, Complexity.prepCondition, Complexity.prepInvariant, Complexity.prepInvariant_run, Complexity.sourcePayload, Complexity.unaryTape, Turing.Cfg, Turing.FinTM, Turing.FinTM.ComputesInTime, Turing.MultiTapeTM.initCfg, Turing.MultiTapeTM.runFrom}
\difficulty{4}
Let $W$ compute on every input $x$ the binary expansion of $T(|x|)$ within
$b\,(T(|x|)+1)$ steps, and suppose $n \le T(n)$ for all $n$.  Then for every
word $x$ there are a time $t$, a base schedule configuration whose unary tape
has $(a+1)(T(|x|)+1)$ markers and whose guide tape has radius
$3(a+1)(T(|x|)+1)$, and data tapes decoding the head-centered payload vectors of
the initial configuration of the source $M$ on $x$, such that the $t$-step run
of the logical simulator on the embedded input is the macrostep configuration
with data at the counter-test phase, unary head $1$, guide head at the origin,
initial data register, and empty output.
\end{lemma}

\begin{lemma}[Iterated macrosteps track the source run]
\lean{Complexity.dataCfg_iterates}
\label{Complexity.dataCfg_iterates}
\leanok
\uses{Complexity.OblPhase, Complexity.OblSymbol, Complexity.blankPayload, Complexity.dataCell, Complexity.dataCfg, Complexity.dataCfg_simulates, Complexity.dataTM, Complexity.guideTape, Complexity.oblEmbed, Complexity.obliviousDataInit, Complexity.sourcePayload, Complexity.sourcePayload_support, Complexity.unaryTape, Complexity.unaryTape_unit, Turing.Cfg, Turing.FinTM, Turing.MultiTapeTM.initCfg, Turing.MultiTapeTM.runFrom, Turing.MultiTapeTM.runFrom_add, Turing.MultiTapeTM.runFrom_succ_eq_step'}
\difficulty{5}
Let the base configuration's unary tape have $B$ markers with $|x| + 1 \le B$
and its guide tape have radius $3B$, and let the data tapes decode the
head-centered payload vectors of the initial configuration of the source $M$ on
$x$.  Then for every $j \le B$ there are reads, neighbor registers, and data
tapes such that $j\,(18B + 8)$ simulator steps from the first counter-test
configuration reach the counter-test phase with unary head $j + 1$, carrying the
state and last-emitted bit of the $j$-step run of $M$ on $x$, with data tapes
decoding that run's head-centered payload vectors and no output emitted.
\end{lemma}

\begin{lemma}[The failed counter test emits the answer]
\lean{Complexity.dataCfg_finish}
\label{Complexity.dataCfg_finish}
\leanok
\uses{Complexity.OblData, Complexity.OblPhase, Complexity.OblSymbol, Complexity.dataCfg, Complexity.dataCfg_step, Complexity.dataTM, Complexity.decoratedCfg, Complexity.macroCfg, Complexity.macroCfg_check, Complexity.macroCfg_unary, Complexity.obliviousVisit, Turing.Cfg, Turing.FinTM, Turing.MultiTapeTM.step}
\difficulty{3}
If the base's unary tape does not hold a unit marker at $u$, then for every data
register and all data tapes, one simulator step from the counter-test phase with
empty output halts, and its output is the single bit symbol carrying the saved
answer bit of the register.
\end{lemma}

\begin{lemma}[The simulator computes the decision bit]
\lean{Complexity.dataTM_computes}
\label{Complexity.dataTM_computes}
\leanok
\uses{Complexity.dataCfg_finish, Complexity.dataCfg_iterates, Complexity.dataTM, Complexity.dataTM_ready, Complexity.oblEmbed, Complexity.sourceAnswer_at_budget, Complexity.unaryTape_end, Turing.FinTM, Turing.FinTM.ComputesInTime, Turing.FinTM.DecidesInTime, Turing.FinTM.computesInTime_iff, Turing.MultiTapeTM.indicator, Turing.MultiTapeTM.initCfg, Turing.MultiTapeTM.runFrom, Turing.MultiTapeTM.runFrom_add, Turing.MultiTapeTM.runFrom_succ_eq_step'}
\difficulty{4}
Let $W$ compute on every input $x$ the binary expansion of $T(|x|)$ within
$b\,(T(|x|)+1)$ steps, suppose $n \le T(n)$ for all $n$, and let $M$ decide the
language $L$ within time $a\,T(n)$.  Then for every word $x$ there is a time $t$
within which the logical simulator, on the embedded input, computes the
one-symbol output consisting of the bit symbol of the indicator of $L$ at $x$.
\end{lemma}

\begin{lemma}[Quadratic decision by the candidate]
\lean{Complexity.obliviousCandidate_decides}
\label{Complexity.obliviousCandidate_decides}
\leanok
\uses{Complexity.dataTM, Complexity.dataTM_computes, Complexity.oblEmbed, Complexity.obliviousCandidate, Complexity.obliviousCandidate_halts, Complexity.parallelTM_computes, Turing.FinTM, Turing.FinTM.ComputesInTime, Turing.FinTM.ComputesInTime.mono, Turing.FinTM.ComputesInTime.output_unique, Turing.FinTM.DecidesInTime, Turing.FinTM.computesInTime_iff, Turing.MultiTapeTM.indicator, Turing.MultiTapeTM.initCfg, Turing.MultiTapeTM.runFrom}
\difficulty{4}
Let $W$ compute on every input $x$ the binary expansion of $T(|x|)$ within
$b\,(T(|x|)+1)$ steps, suppose $n \le T(n)$ for all $n$, and let $M$ decide the
language $L$ within time $a\,T(n)$.  Then the candidate simulator built from
$W$, $M$, and $a$ decides $L$ within time
$\bigl(18(a+1)^2 + 23(a+1) + 3b + 25\bigr)\,(T(n)+1)^2$.
\end{lemma}

\begin{theorem}[Oblivious simulation with quadratic overhead]
\lean{Complexity.oblivious_of_mem_DTIME}
\label{Complexity.oblivious_of_mem_DTIME}
\leanok
\uses{Complexity.DTIME, Complexity.TimeConstructible, Complexity.obliviousCandidate, Complexity.obliviousCandidate_decides, Complexity.obliviousCandidate_oblivious, Turing.FinTM, Turing.FinTM.DecidesInTime, Turing.FinTM.Oblivious}
\difficulty{2}
Let $T$ be a time-constructible bound and let $L$ be a binary language in
$\mathrm{DTIME}(T)$.  Then there exist an oblivious machine $M$ — one whose head
positions at every time step depend only on the input length — and a constant
$c$ such that $M$ decides $L$ within time $c\,(T(n)+1)^2$.  This is the first
assertion of Arora--Barak, Exercise~1.5.
\end{theorem}
